\begin{table}[htbp]
\centering
\caption{Minimal winning coalitions and membership by ideological domain}
\label{tab:interval-summary}
\small
\setlength{\tabcolsep}{4pt}
\begin{tabularx}{\textwidth}{@{}rc>{\centering\arraybackslash}Xc cc@{}}
\toprule
 & & & & \multicolumn{2}{c}{\shortstack{Coalition members:\\median [min--max]}} \\
\cmidrule(l){5-6}
Election & \(k\) & Minimal majorities & Inversions & Inverted & Non-inverted \\
\midrule
2014 & 0 & 8 & 4 & 15 [12--16] & 12 [12--14] \\
2014 & 1 & 88 & 43 & 14 [11--19] & 14 [11--19] \\
2018 & 0 & 8 & 0 & None & 16.5 [15--20] \\
2018 & 1 & 118 & 24 & 17 [14--19] & 17 [14--21] \\
2022 & 0 & 4 & 2 & 11.5 [8--15] & 20 [19--21] \\
2022 & 1 & 53 & 17 & 14 [7--16] & 19 [14--21] \\
\bottomrule
\end{tabularx}
\begin{flushleft}
\footnotesize Notes: Coalition membership is the number of parties in the coalition.
Entries in the final two columns report median [minimum--maximum].
Membership counts include parties receiving votes but no seats; they are not counts of seat-winning parliamentary partners.
Minimality is defined within the stated domain.
\end{flushleft}
\end{table}
